\section{Discussion}
DIReCT$++$ combines shared tracer-conditioned 3D rectified flow with biomedical text guidance for FDG- and AV45-PET synthesis from MRI. Its contribution extends the single-tracer DIReCT framework to multi-tracer modeling and a separately evaluated affine-adapter formulation. The evidence comprises historical aggregate comparisons, reconstructed OOF generation, normalized regional uptake statistics and a separate cross-sectional classification analysis. These sources should be interpreted within their respective input and execution identities rather than as a single verified historical run.

The generation results show both advantages and tradeoffs. Full had better masked image-quality means than the clinical regression U-Net and MRI-only Swin-UNETR, while some AV45 error effects were small and did not imply universal individual benefit. Whole-volume comparisons favored TabularMLP for PSNR/MSE and TracerOnly for several endpoints. The adapter analysis supports control of tracer-embedding geometry, but it does not establish a causal improvement in PET generation. Fixed-weight conditioning interventions indicate that patient features can influence outputs; independently initialized Full/TracerOnly or text/tabular models do not isolate the clinical or VLM component as the sole cause of their performance differences.

The potential value of this framework is to improve access to biomarker-related imaging information for multimodal imaging research. The classification findings are task dependent: MRI+Syn-both improved AD/CN relative to MRI, did not improve MCI/CN, and showed mixed EMCI/LMCI endpoints. Clinical-only and MRI+Clinical AUC exceeded MRI+Syn-both in the completed comparisons. Because clinical inputs include diagnosis-related assessments and diagnosis-conditioned fills, these findings cannot be used to establish independent diagnostic value of synthetic PET. Similar observed Real/Syn performance does not demonstrate equivalence or molecular validity.

Synthetic PET may support exploratory MRI-based analyses when PET is unavailable, but clinical screening or treatment selection would require independent molecular biomarker confirmation and prospective validation. No longitudinal conversion, progressive-versus-stable MCI, survival or treatment-response endpoints were analyzed. EMCI/LMCI classification is therefore MCI stage discrimination; it does not predict an individual's progression to AD. Whether serially synthesized PET can reflect progression or treatment response remains a question for longitudinal studies.

Within medical image synthesis, DIReCT$++$ evaluates disease-related regional uptake patterns and downstream classification in addition to image-level metrics. The shared flow model and biomedical conditioning provide a framework for further study, rather than establishing clinically actionable biomarkers or attributing all overall differences to one architectural component. The learned one-step approximation permits a single tracer-conditioned forward pass; measured forward latency does not constitute an end-to-end workflow or an equal-protocol speedup over competing models.

Several limitations require explicit consideration. Clinical-record reconstruction showed extensive diagnosis-conditioned filling of missing scores, which can introduce group information into conditioning and classification. Source weight entries also contain unresolved validity issues. The original random filling states and complete actual acquisition mappings were not recovered. The independent three-fold classifier evaluation verifies its own partitions, but the historical synthetic-image inputs could not be bound directly to the checkpoints in the reconstructed generation comparisons. These limitations constrain claims of reproducibility and independent clinical discrimination.

Regional measurements use per-image min--max-normalized intensities in arbitrary units, without established anatomical-reference SUV/SUVR calibration. No clinically justified equivalence margin, formal TOST result or valid amyloid-positivity agreement analysis was established. Close agreement, non-significant paired differences or preservation of AD--CN significance is not evidence of clinical equivalence or reliable atypical A$\beta$ status classification. Calibrated quantitative PET and prespecified, clinically meaningful equivalence margins are needed for such assessments.

The deterministic one-step formulation provides a single output per input without calibrated subject-level predictive uncertainty. An MSE-trained deterministic model may represent an average conditional mapping and cannot guarantee recovery of individual molecular pathology from MRI and text. Ensemble or stochastic uncertainty estimation remains future work. Not all requested single-component ablations or matched one-step/multistep evaluations were completed. The fixed-model bootstrap intervals do not quantify training-seed or split variability.

External evaluation was limited to historical AV45 synthesis in OASIS-3 CN/AD, with age differences and incomplete run-specific prompt records; FDG and MCI external generalizability remain unestablished. Prospective validation in more diverse clinical populations, including alternative MRI acquisitions~\cite{sorby2024portable}, is needed. Additional clinical modalities such as blood biomarkers~\cite{hansson2023blood} may provide complementary information, but their incremental value must be evaluated independently. The 2024 Alzheimer's Association criteria emphasize biological diagnosis using validated core biomarkers, including amyloid PET and sufficiently accurate blood-based assays such as plasma p-tau217~\cite{jack2024revised}. Synthetic PET has not been validated as a substitute for molecular confirmation, and its added value relative to those biomarker assessments remains to be determined.

In conclusion, DIReCT$++$ provides a shared tracer-conditioned framework for PET synthesis with biomedical text inputs and learned one-step inference. The reported image-quality and regional results support further methodological investigation, while mixed classification effects, clinical-input circularity and unresolved source limitations constrain clinical interpretation. Longitudinal prognosis, clinical equivalence and replacement of molecular biomarker confirmation require separate validation.
